\chapter{Conclusões e Trabalhos Futuros}\label{cap6-conclusoes}

Este capítulo sintetiza as conclusões frente aos objetivos formulados na \secref{sec-13-objetivo-geral}, destaca as principais contribuições do trabalho e propõe direções para trabalhos futuros.

\section{Conclusões}\label{sec-61-conclusoes}

Como arquiteturas de CPU e GPU possuem modelos de execução distintos, um algoritmo não pode ser implementado de maneira estritamente idêntica nas duas plataformas: cada variante exige otimizações próprias (como vetorização na CPU e distribuição de \emph{threads}, uso de memória compartilhada e operações com máscaras de bits na GPU).
Para tornar a comparação justa dentro desse limite real, o estudo manteve comuns o núcleo do algoritmo memético AG+\(2\)-opt, os parâmetros globais e os critérios de parada, enquadrando as estratégias aceleradas no paralelismo de baixo nível (Tipo~1) da taxonomia de \textcite{crainic2003parallel,crainic2010parallel} (\secref{sec-261-taxonomia-paralelizacao}) e avaliando as quatro variantes em \totalinstancias\ instâncias da TSPLIB (de \(\menorinstancia\) a \(\maiorinstancia\) cidades, nos formatos \texttt{EUC\_2D}, \texttt{ATT}, \texttt{GEO} e \texttt{EXPLICIT}) com \totalrepeticoes\ repetições pareadas (\capref{cap3-materiais-metodos}).

Os resultados do \capref{cap4-resultados} mostram que a estratégia de paralelização determina fortemente o desempenho sem comprometer a qualidade das soluções.
Na métrica de qualidade (\emph{gap}), o protocolo estatístico de \textcite{demsar2006statistical} não identificou diferença estatisticamente significativa entre as quatro variantes (teste de Friedman, \(p = 0{,}4543\), com \emph{gaps} médios entre \(0{,}90\%\) e \(0{,}91\%\), e \(0\) de \(90\) pares significativos no Eixo~1, \secref{sec-341-eixo1}).
Já no tempo de execução, a variante \textbf{Totalmente em GPU} obteve o menor tempo em todas as \totalinstancias\ instâncias (ganho médio geométrico de \(71{,}03\times\) sobre a CPU e tempo mediano de \(0{,}34\)\,s), seguida pela \textbf{Híbrida Otimizada} (\(23{,}86\times\); mediana de \(1{,}02\)\,s) e pela \textbf{Híbrida Ingênua} (\(4{,}57\times\); mediana de \(4{,}81\)\,s), com superioridade estatisticamente significativa nos seis pares em todas as \totalinstancias\ instâncias (\(p_{\text{Holm}} < 0{,}001\)).

A análise por etapa (\secref{sec-44-decomposicao-fases}) confirma ainda que reduzir o tráfego CPU--GPU é condição necessária, mas não suficiente, para extrair desempenho do acelerador: o ganho prático resulta da combinação entre baixa comunicação, \emph{kernels} bem estruturados e equilíbrio do ciclo evolutivo frente à Lei de Amdahl (Equação~\eqref{eq:amdahl}, como visto na inversão de gargalo para o cruzamento em CPU na Híbrida Otimizada).
Por fim, embora o limite de VRAM não tenha sido atingido no conjunto avaliado, a literatura aponta que o crescimento quadrático da matriz densa tende a se tornar restritivo em problemas de porte muito maior \autocite{schulz2013gpu,oneil2015rethinking}, exigindo alternativas como cálculo sob demanda, fatiamento em memória compartilhada ou estratégias híbridas, ainda que em geral com menor eficiência temporal.

\section{Contribuições}\label{sec-62-contribuicoes}

As principais contribuições deste trabalho podem ser reunidas em quatro frentes:

\begin{itemize}
      \item
            \textbf{Contribuição empírica:} uma comparação controlada de quatro estratégias de paralelização sobre o mesmo algoritmo memético (AG+\(2\)-opt), mantendo os mesmos parâmetros e critérios de parada para avaliar o impacto prático da organização dos dados e dos \emph{kernels} em GPU;
      \item
            \textbf{Contribuição metodológica:} um protocolo experimental reproduzível, com múltiplas repetições, sementes fixas e testes estatísticos não paramétricos adequados a comparações múltiplas;
      \item
            \textbf{Contribuição prática (ferramenta experimental):} um \emph{framework} modular em três camadas (algoritmos, orquestração e análise) que permite alternar implementações de CPU e GPU sob a mesma interface, com medição por etapa, contagem de transferências e registro estruturado dos experimentos; e
      \item
            \textbf{Contribuição analítica:} a caracterização do compromisso entre tempo, qualidade e comunicação em uma GPU de entrada, mostrando como os gargalos se deslocam entre busca local, transferências PCIe e cruzamento em CPU conforme a Lei de Amdahl.
\end{itemize}

\section{Trabalhos futuros}\label{sec-64-trabalhos-futuros}

Como desdobramento deste estudo, destacam-se cinco frentes de investigação:

\begin{itemize}
      \item
            \textbf{Outros operadores genéticos e de busca local em GPU:} Este trabalho fixou operadores clássicos (OX, mutação por troca e busca local \(2\)-opt) para manter a comparação controlada. Um passo natural é investigar como operadores mais sofisticados de busca local e recombinação --- como \(3\)-opt, movimentos inspirados em Lin--Kernighan ou cruzamentos baseados em arestas, como o EAX \autocite{nagata2013powerful} --- se comportam quando adaptados para GPU, equilibrando ganho de qualidade e custo de paralelização.
      \item
            \textbf{Escalonamento para instâncias de grande porte e limites de memória:} Nas instâncias avaliadas (até \(\maiorinstancia\) cidades), o consumo de VRAM permaneceu muito abaixo da capacidade da placa. Resta avaliar empiricamente o desempenho quando o porte do problema ultrapassa os limites de memória compartilhada ou da VRAM, comparando estratégias como cálculo sob demanda, particionamento de dados e abordagens híbridas em instâncias com dezenas de milhares de cidades \autocite{schulz2013gpu,oneil2015rethinking}.
      \item
            \textbf{Extensão para variantes mais realistas de roteamento:} Estender a investigação para problemas mais próximos de aplicações industriais e logísticas --- como o TSP assimétrico (ATSP), o Problema de Roteamento de Veículos com Capacidade (CVRP) ou problemas com janelas de tempo --- permitirá verificar como restrições adicionais e desvios condicionais afetam o aproveitamento da GPU.
      \item
            \textbf{Generalização entre arquiteturas de \emph{hardware} e custo energético:} Repetir o protocolo em diferentes combinações de CPU (incluindo versões paralelizadas em múltiplos núcleos) e GPUs de gerações distintas, incorporando também a medição de consumo de energia, ajudará a delimitar quando o uso de GPU oferece melhor relação entre desempenho e custo computacional.
      \item
            \textbf{Comparação estatística sob abordagem bayesiana:} A análise conduzida neste trabalho é frequentista, verificando se há evidência para rejeitar a hipótese de equivalência sem estimar diretamente a probabilidade de uma variante superar outra. Como propuseram \textcite{benavoli2017time} em revisão ao protocolo de 2006, adotar testes bayesianos sobre os dados pareados permitiria substituir o veredito binário por probabilidades diretas de superioridade temporal ou de equivalência prática em qualidade, mediante a definição prévia de uma margem de equivalência prática e de uma distribuição inicial.
\end{itemize}
